\documentclass[11pt,a4paper]{article}
\usepackage{fontspec}
\setmainfont{Malgun Gothic}
\usepackage[margin=2cm]{geometry}
\usepackage{graphicx}
\usepackage{xcolor}
\usepackage{titlesec}
\usepackage{enumitem}
\usepackage{tabularx}
\usepackage{booktabs}
\usepackage{fancyhdr}
\usepackage{hyperref}
\usepackage{float}
\usepackage{needspace}
\usepackage{multicol}

\definecolor{korfarm}{HTML}{F06C24}
\definecolor{korfarmDark}{HTML}{2B221D}
\definecolor{korfarmGreen}{HTML}{4E8A60}
\definecolor{korfarmBlue}{HTML}{4D8BCC}
\definecolor{lightbg}{HTML}{FFF4E9}

\titleformat{\section}{\Large\bfseries\color{korfarm}}{}{0em}{}[\titlerule]
\titleformat{\subsection}{\large\bfseries\color{korfarmDark}}{}{0em}{}
\titleformat{\subsubsection}{\normalsize\bfseries\color{korfarmGreen}}{}{0em}{}

\pagestyle{fancy}
\fancyhf{}
\fancyhead[L]{\small\color{gray}국어농장 학생 학습 안내서}
\fancyhead[R]{\small\color{gray}\thepage / \pageref{LastPage}}
\fancyfoot[C]{\small\color{gray}국어농장 --- 스스로 완성하는 국어 근육}

\usepackage{lastpage}

\hypersetup{colorlinks=true,linkcolor=korfarm,urlcolor=korfarmBlue}

\newcommand{\screenshot}[2]{%
  \begin{figure}[H]
  \centering
  \fbox{\includegraphics[width=#2\textwidth,keepaspectratio]{screenshots/#1}}
  \end{figure}
}

\newcommand{\goalbox}[1]{%
  \vspace{0.3cm}
  \noindent\colorbox{lightbg}{\parbox{\dimexpr\textwidth-2\fboxsep}{%
    \textbf{\color{korfarm}학습 목표} \quad #1%
  }}
  \vspace{0.3cm}
}

\begin{document}

%% ═══════════════════════════════════════════
%% 표지
%% ═══════════════════════════════════════════
\begin{titlepage}
\centering
\vspace*{3cm}
{\Huge\bfseries\color{korfarm} 국어농장\par}
\vspace{0.3cm}
{\LARGE 스스로 완성하는 국어 근육\par}
\vspace{1.5cm}
{\huge\bfseries 학생 학습 기능 안내서\par}
\vspace{0.8cm}
{\large 초등학교 1학년 \textasciitilde{} 고등학교 3학년 대상\par}
\vspace{3cm}

\begin{tabular}{rl}
\textbf{레벨 체계} & 소쉬르(초1\textasciitilde{}3) · 프레게(초4\textasciitilde{}6) · 러셀(중1\textasciitilde{}3) · 비트겐슈타인(고1\textasciitilde{}3) \\
\textbf{10대 핵심 역량} & 어휘력 · 독해력 · 문법능력 · 논리사고력 · 배경지식 외 \\
\textbf{학습 영역} & 독해 · 어휘 · 문법 · 배경지식 · 논리사고력 · 글쓰기 \\
\end{tabular}

\vfill
{\large 2026년 4월\par}
\end{titlepage}

\tableofcontents
\newpage

%% ═══════════════════════════════════════════
\section{학습 시작 화면}
%% ═══════════════════════════════════════════

로그인 후 처음 만나는 나만의 학습 홈 화면입니다.

\screenshot{01_start.png}{0.9}

\subsection{화면 구성}
\begin{itemize}[nosep,leftmargin=*]
  \item \textbf{상단 프로필}: 이름, 현재 레벨, 시즌 점수가 표시됩니다. 내 정보 수정, 구독 관리, 로그아웃 버튼이 있습니다.
  \item \textbf{검색창}: 학습 콘텐츠를 키워드로 검색할 수 있습니다. 일일퀴즈, 일일독해, 프로모드, 농장모드 카테고리로 필터링됩니다.
  \item \textbf{오늘의 무료 학습}: 일일 퀴즈와 일일 독해 카드. 현재 레벨과 진행 일차가 표시됩니다.
  \item \textbf{과제 바구니}: 선생님이 배정한 특별 과제를 확인합니다.
  \item \textbf{학습 메뉴}: 프로 모드, 농장별 모드, 지식과 지혜, 테스트 창고, 역량 진단, 시험 공부.
  \item \textbf{대결 모드}: 실시간 대결 신청과 대결 랭킹 TOP 5.
  \item \textbf{통계·기록}: 시즌 랭킹, 테스트 기록실, 통합 성적표.
  \item \textbf{하단}: 수확 장부, 씨앗 획득 내역, 씨앗 교환, 커뮤니티, 쇼핑몰.
\end{itemize}

\newpage
%% ═══════════════════════════════════════════
\section{일일 학습 (무료)}
%% ═══════════════════════════════════════════

매일 무료로 제공되는 학습입니다. 반복 학습이 가능하며, 각각 하루 씨앗 10개까지 획득합니다.

\subsection{일일 퀴즈}
\goalbox{레벨에 맞는 국어 문제를 매일 풀며 국어 실력의 기초를 다집니다.}
\begin{itemize}[nosep,leftmargin=*]
  \item 매일 학생의 레벨에 맞는 \textbf{10문제}가 자동 출제됩니다.
  \item 제한 시간이 화면 상단에 타이머로 표시됩니다.
  \item \textbf{정답}이면 시간이 추가(+20초)되고, \textbf{오답}이면 시간이 감소(-20초)합니다.
  \item 시험지 형식으로 문제 목록이 좌우 2단으로 배치됩니다.
  \item 예문이 있는 문제는 예문 박스에 밑줄 친 단어가 하이라이트됩니다.
  \item 오답이면 사선 표시 후 다음 문제로 자동 진행됩니다.
  \item 총 10문제 완료 시 결과가 표시되고 씨앗이 지급됩니다.
\end{itemize}

\subsection{일일 독해}
\goalbox{지문을 정독하고, 핵심을 복기하며, 정답을 직접 찾는 3단계 독해 훈련입니다.}

독해 학습은 \textbf{3단계}로 진행됩니다:

\subsubsection{1단계: 정독}
\begin{itemize}[nosep,leftmargin=*]
  \item 지문이 표시되고, 타임라인에 따라 핵심 표현이 순차적으로 \textbf{하이라이트}됩니다.
  \item 하이라이트된 부분에 대한 \textbf{질문}이 모달로 출제됩니다.
  \item 정답이면 시간이 추가되고, 오답이면 감소합니다.
  \item 타임라인의 모든 질문을 완료하면 복기 단계로 넘어갑니다.
\end{itemize}

\subsubsection{2단계: 복기}
\begin{itemize}[nosep,leftmargin=*]
  \item 지문의 핵심 내용이 \textbf{카드} 형식으로 표시됩니다.
  \item 각 카드에서 빠진 내용을 맞춥니다.
  \item 틀리면 씨앗 페널티가 적용됩니다.
\end{itemize}

\subsubsection{3단계: 확인}
\begin{itemize}[nosep,leftmargin=*]
  \item ``\textbf{다음 질문의 답을 지문에서 찾아 클릭하세요}''라는 안내가 표시됩니다.
  \item 지문이 글자 단위로 클릭 가능한 상태가 됩니다.
  \item 질문을 읽고, 지문에서 해당하는 부분을 \textbf{직접 클릭}합니다.
  \item 정답 범위를 정확히 클릭하면 정답 처리되고, 다음 확인 문제로 넘어갑니다.
  \item ALL 모드인 경우 여러 범위를 모두 찾아야 합니다.
\end{itemize}

\newpage
%% ═══════════════════════════════════════════
\section{농장별 모드}
%% ═══════════════════════════════════════════
\goalbox{11개 주제별 농장에서 원하는 영역을 골라 집중 학습합니다. 각 농장은 특정 씨앗을 보상으로 지급합니다.}

\screenshot{02_farm_mode.png}{0.85}

\subsection{농장 목록과 대상 레벨}

\begin{tabularx}{\textwidth}{llXl}
\toprule
\textbf{이모지} & \textbf{농장} & \textbf{설명} & \textbf{대상} \\
\midrule
🌾 & 어휘 농장 & 낱말의 뜻, 용례, 사전 활용 학습 & 전체 \\
📖 & 독해 농장 & 비문학·문학 지문을 읽고 독해 훈련 & 전체 \\
🌍 & 배경지식 농장 & 다양한 주제의 배경지식 퀴즈 & 전체 \\
📚 & 이야기 농장 & 소설·수필·동화 문학 독해 & 초등 \\
📜 & 고전 농장 & 고전 시가·산문 독해 & 초등 \\
📋 & 내용 숙지 농장 & PDF/텍스트 기반 확인 학습 & 중·고등 \\
🔤 & 문법 농장 & 형태소, 단어형성, 문장구조, 음운변동, 품사 & 초고\textasciitilde{}고등 \\
💡 & 국어 개념 농장 & 국어 핵심 개념과 이론 & 중·고등 \\
🧩 & 논리사고력 농장 & 추론과 논리적 사고 훈련 & 초고\textasciitilde{}고등 \\
✍️ & 서술형 농장 & 서술형 문제 풀이 연습 & 중·고등 \\
✅ & 선택지 판별 농장 & 선택지 적절성 판별 훈련 & 고등 \\
\bottomrule
\end{tabularx}

\subsection{학습 흐름}
\begin{enumerate}[nosep,leftmargin=*]
  \item 농장별 모드 화면에서 레벨(서버)를 선택합니다.
  \item 접근 가능한 농장이 카드 형태로 표시되고, 각 농장의 학습 콘텐츠 개수가 표시됩니다.
  \item 농장을 클릭하면 해당 농장의 세부 학습 목록이 표시됩니다.
  \item 학습을 선택하면 학습 엔진이 실행됩니다.
  \item 학습 완료 시 씨앗이 지급되고, 수확 장부에 기록됩니다.
\end{enumerate}

\newpage
%% ═══════════════════════════════════════════
\section{문법 학습 상세}
%% ═══════════════════════════════════════════
\goalbox{한국어 문법의 5대 영역(품사, 형태소, 단어형성, 문장짜임, 음운변동)을 단계별로 학습합니다.}

\screenshot{15_farm_grammar.png}{0.85}

\subsection{품사 학습 (기초 50개 + 심화 50개, 총 2,000문항)}
\goalbox{예문 속 단어의 품사를 정확하게 판별하는 능력을 기릅니다. 표준국어대사전 기반 예문.}

\subsubsection{기초 (프레게+러셀 대상)}
\begin{itemize}[nosep,leftmargin=*]
  \item 품사 통용이 없는 명확한 단어를 대상으로 합니다.
  \item 예: ``\textbf{아름다운} 꽃이 피었다.'' → 형용사
  \item 명사, 동사, 형용사, 부사, 관형사, 감탄사, 조사, 대명사, 수사를 판별합니다.
\end{itemize}

\subsubsection{심화 (러셀+비트겐슈타인 대상)}
\begin{itemize}[nosep,leftmargin=*]
  \item 문맥에 따라 품사가 달라지는 \textbf{품사 통용} 단어를 대상으로 합니다.
  \item 예: ``나\textbf{보다} 크다'' → 조사 / ``영화를 \textbf{보다}'' → 동사
  \item 보조 동사, 보조 형용사, 의존 명사 등 세부 품사도 포함합니다.
\end{itemize}

\subsubsection{학습 방법}
\begin{enumerate}[nosep,leftmargin=*]
  \item 시험지 형식으로 20문제가 좌우 2단으로 배치됩니다.
  \item 현재 풀어야 할 문제가 하이라이트됩니다.
  \item 문제 아래에 모달이 뜨고, 4지선다로 품사를 선택합니다.
  \item 예문 박스 안의 밑줄 친 단어가 출제 대상입니다.
  \item 오답이면 사선 표시 후 다음 문제로 자동 진행됩니다.
\end{enumerate}

\newpage
\subsection{형태소 분석 (기초 10개 + 심화 10개)}
\goalbox{문장을 형태소 단위로 분해하고, 각 형태소의 이름(명사, 조사, 어간, 어미 등)과 종류(실질 자립/실질 의존/형식 의존)를 판별합니다.}

\subsubsection{학습 방법 (4단계)}
\begin{enumerate}[nosep,leftmargin=*]
  \item \textbf{1단계 — 형태소 개수}: 문장 전체가 하이라이트되고, ``이 문장의 형태소는 몇 개인가요?'' 4지선다.
    \begin{itemize}[nosep]
      \item 예: ``나는 학교에 간다.'' → 정답: 7개 (나, 는, 학교, 에, 가-, -ㄴ-, -다)
    \end{itemize}
  \item \textbf{2단계 — 형태소 구분}: 올바르게 구분한 선택지를 고릅니다. 정답이면 문장이 분해된 형태소 칩으로 교체됩니다.
    \begin{itemize}[nosep]
      \item 정답 예: ``나, 는, 학교, 에, 가-, -ㄴ-, -다''
      \item 오답 예: ``나는, 학교에, 간다'' (어절 단위 — 미분리)
    \end{itemize}
  \item \textbf{3단계 — 형태소 이름}: 각 형태소 칩이 순차적으로 하이라이트되며, 이름을 4지선다로 맞춥니다.
    \begin{itemize}[nosep]
      \item 기초: 명사, 대명사, 조사, 어간, 어미, 접사 등 큰 분류
      \item 심화: 3인칭 대명사, 보조사, 동사 어간, 선어말 어미, 사동 접미사 등 세부 분류
    \end{itemize}
  \item \textbf{4단계 — 형태소 종류}: 같은 형태소에 대해 3지선다로 종류를 맞춥니다.
    \begin{itemize}[nosep]
      \item \textbf{실질 자립}: 명사, 대명사, 수사, 관형사, 부사, 감탄사
      \item \textbf{실질 의존}: 동사/형용사 어간, 비자립성 어근
      \item \textbf{형식 의존}: 조사, 어미, 접사
    \end{itemize}
\end{enumerate}

\subsubsection{형태소 분석 규칙}
\begin{itemize}[nosep,leftmargin=*]
  \item 준말은 원형으로 풀어서 분석합니다. (왔다 → 오- + -았- + -다)
  \item 어간/어미/접사는 기본형으로 표시합니다. (서둘러 → 서두르- + -어)
  \item ``-ㄴ다''는 ``-ㄴ-''(선어말 어미) + ``-다''(종결 어미)로 분리합니다.
  \item 서술격 조사 ``이다'' → ``이-''(조사) + ``-다''(어미)
  \item 사동/피동 접미사는 별도 분리합니다. (밝히다 → 밝- + -히- + -다)
\end{itemize}

\newpage
\subsection{단어의 형성 (기초 20개 + 심화 10개)}
\goalbox{단어가 어떤 형태소의 결합으로 만들어졌는지 분석합니다 (단일어, 합성어, 파생어).}

\subsubsection{학습 방법 (5단계)}
\begin{enumerate}[nosep,leftmargin=*]
  \item \textbf{1\textasciitilde{}4단계}: 형태소 분석과 동일한 4단계를 거칩니다.
  \item 4단계 완료 후, 각 형태소에 표시가 됩니다:
    \begin{itemize}[nosep]
      \item {\color{korfarmGreen}\textbf{동그라미}}: 어근/실질 형태소 (녹색 원형 테두리)
      \item {\color{korfarm}\textbf{네모}}: 접사 (주황색 사각 테두리)
      \item 어미는 표시 없음
    \end{itemize}
  \item \textbf{5단계}: ``이 단어는 어떤 단어인가요?'' 3지선다.
    \begin{itemize}[nosep]
      \item \textbf{단일어}: 하나의 어근 (예: 구름, 하늘, 나무)
      \item \textbf{합성어}: 어근 + 어근 (예: 눈사람, 밤낮, 빛나다)
      \item \textbf{파생어}: 접사 + 어근 또는 어근 + 접사 (예: 맨손, 먹이, 헛기침)
    \end{itemize}
\end{enumerate}

\subsubsection{심화 추가 단계}
\begin{itemize}[nosep,leftmargin=*]
  \item 합성어인 경우: \textbf{통사적 합성어}인지 \textbf{비통사적 합성어}인지 판별.
  \item 합성어 관계: \textbf{대등}, \textbf{종속}, \textbf{융합} 중 선택.
  \item 3개 이상 형태소인 경우: 결합 순서를 단계적으로 분석합니다.
    \begin{itemize}[nosep]
      \item 예: ``거짓말쟁이'' → 거짓+말(합성어) → 거짓말+-쟁이(파생어)
    \end{itemize}
\end{itemize}

\newpage
\subsection{문장의 짜임 (기초 20개 + 심화 15개)}
\goalbox{문장의 구조를 분석하고, 겹문장의 절 유형(이어진 문장, 안긴 문장)을 판별합니다.}

\subsubsection{학습 방법}
\begin{enumerate}[nosep,leftmargin=*]
  \item 문장이 \textbf{어절 단위 박스}로 표시됩니다. (본용언+보조용언은 한 박스)
  \item \textbf{서술어부터} 시작하여 각 어절을 클릭합니다.
  \item 클릭하면 모달이 뜨고, 문장성분을 5지선다(레이어 포함)로 선택합니다.
    \begin{itemize}[nosep]
      \item 예: ``주1'', ``서2'', ``목1'', ``부1'' 등
    \end{itemize}
  \item 같은 레이어의 성분을 모두 채운 후, 다음 레이어의 서술어를 클릭합니다.
  \item 모든 문장성분을 채우면 \textbf{겹문장 분석}으로 넘어갑니다.
  \item 절의 \textbf{시작 칸과 마지막 칸}을 클릭하여 범위를 선택합니다.
  \item 선택한 범위가 어떤 역할인지 판별합니다:
    \begin{itemize}[nosep]
      \item \textbf{이어진 문장}: 절 경계에 사선(/) 표시 → 대등/종속 판별
      \item \textbf{안긴 문장}: 괄호 표시 + 행 추가 → 명사절/관형절/부사절/서술절/인용절 판별
    \end{itemize}
\end{enumerate}

\subsubsection{기초 예시}
``그가 범인임이 확실하다.''
\begin{itemize}[nosep,leftmargin=*]
  \item 1열: 주2 | 서2 | 서1
  \item 2열: (주1 명사절) — ``그가 범인임이'' 범위
\end{itemize}

\subsubsection{심화 예시}
``나는 그가 준 책을 잃어버린 기억을 잊고 싶다.''
\begin{itemize}[nosep,leftmargin=*]
  \item 1열: 주1 | 주3 | 서3 | 목2 | 서2 | 목1 | 서1
  \item 2열: (관2 관형절) — ``그가 준'' 범위
  \item 3열: (관1 관형절) — ``그가 준 책을 잃어버린'' 범위
\end{itemize}

\newpage
\subsection{음운 변동 (러셀+비트겐슈타인 대상)}
\goalbox{단어의 기저형(시작점)에서 표면형(도착점)으로 바뀌는 음운 변동 과정을 분석합니다.}

\subsubsection{학습 방법}
\begin{enumerate}[nosep,leftmargin=*]
  \item 단어의 \textbf{시작점}(기저형)이 음운 칸으로 표시됩니다.
  \item 아래에 \textbf{도착점}(표면형) 칸이 표시됩니다.
  \item 변동이 일어나는 칸을 클릭합니다. 모든 칸이 클릭 가능합니다.
  \item 정답 칸을 클릭하면 모달이 뜨고, \textbf{변동 결과 음운}을 선택합니다.
  \item 그 후 적용된 \textbf{음운 변동 규칙}을 맞춥니다.
  \item 틀린 칸을 클릭하면 \textbf{흔들림 애니메이션}과 함께 오답 처리(-40초)됩니다.
  \item 모든 칸은 동일한 정사각형(44×44px)으로 표시됩니다.
\end{enumerate}

\subsubsection{음운 변동 예시}
``국물'' [궁물]
\begin{itemize}[nosep,leftmargin=*]
  \item 시작: ㄱ | ㅜ | ㄱ | , | ㅁ | ㅜ | ㄹ
  \item 도착: ㄱ | ㅜ | ㅇ | , | ㅁ | ㅜ | ㄹ (ㄱ→ㅇ 비음화)
\end{itemize}

\newpage
%% ═══════════════════════════════════════════
\section{프로 모드}
%% ═══════════════════════════════════════════
\goalbox{12레벨 순차 심화 학습. 이전 챕터를 통과해야 다음 챕터가 해제되는 체계적 학습 시스템입니다.}

\screenshot{03_pro_mode.png}{0.85}

\begin{itemize}[nosep,leftmargin=*]
  \item 12레벨 × 다수 챕터로 구성됩니다.
  \item 각 챕터에는 다양한 학습이 포함됩니다: \textbf{독해}, \textbf{어휘}, \textbf{배경지식}, \textbf{논리사고력}, \textbf{모범답안/정답해설}.
  \item 챕터별 \textbf{진행률(\%)}이 표시됩니다.
  \item \textbf{영상 학습 링크}가 제공됩니다.
  \item 잠금 해제: 이전 챕터의 테스트를 통과해야 합니다.
  \item 구독 회원 전용입니다.
\end{itemize}

\newpage
%% ═══════════════════════════════════════════
\section{지식과 지혜 (글쓰기)}
%% ═══════════════════════════════════════════
\goalbox{레벨에 맞는 주제로 글을 쓰고 선생님의 첨삭을 받습니다. 사고력과 표현력을 함께 기릅니다.}

\screenshot{07_writing.png}{0.85}

\subsection{레벨별 글쓰기 주제}
\begin{tabularx}{\textwidth}{lX}
\toprule
\textbf{레벨} & \textbf{주제 예시} \\
\midrule
소쉬르 1\textasciitilde{}3 (초1\textasciitilde{}3) & 계절, 가족, 감정 표현 등 생활 밀착형 글쓰기 \\
프레게 1\textasciitilde{}3 (초4\textasciitilde{}6) & 과학 탐구, 경제 기초, 독서 감상 글쓰기 \\
러셀 1\textasciitilde{}3 (중1\textasciitilde{}3) & 고전 철학, 헌법 원리, 과학 법칙 글쓰기 \\
비트겐슈타인 1\textasciitilde{}3 (고1\textasciitilde{}3) & 인식론, 금융 시장, 양자역학, 생명윤리 글쓰기 \\
\bottomrule
\end{tabularx}

\subsection{학습 흐름}
\begin{enumerate}[nosep,leftmargin=*]
  \item 12개 레벨 중 자신의 레벨을 선택합니다.
  \item 주어진 주제에 대해 자유롭게 글을 작성합니다.
  \item 제출하면 선생님이 \textbf{첨삭}합니다.
  \item 첨삭 결과를 확인하고, 개선된 글을 다시 작성할 수 있습니다.
\end{enumerate}

\newpage
%% ═══════════════════════════════════════════
\section{대결}
%% ═══════════════════════════════════════════
\goalbox{다른 학생과 실시간으로 국어 문제 대결을 합니다. 씨앗을 걸고 실력을 겨루세요!}

\screenshot{05_duel.png}{0.85}

\subsection{서버 선택}
수준에 맞는 서버를 선택합니다:
\begin{itemize}[nosep,leftmargin=*]
  \item 🌱 \textbf{소쉬르} — 초등 저학년 수준
  \item 🌾 \textbf{프레게} — 초등 고학년 수준
  \item 🌿 \textbf{러셀} — 중학교 수준
  \item 🌳 \textbf{비트겐슈타인} — 고등학교 수준
\end{itemize}

\subsection{대결 진행}
\begin{enumerate}[nosep,leftmargin=*]
  \item 서버의 로비에서 열린 방에 입장하거나, 새 방을 만듭니다.
  \item 방을 만들 때 이름과 \textbf{베팅 씨앗}(1\textasciitilde{}50개)을 설정합니다.
  \item 참가자가 모이면 대결이 시작됩니다.
  \item 매 라운드 \textbf{제한 시간}(30초) 안에 문제를 풉니다.
  \item 오답자는 \textbf{탈락}하고, 마지막까지 남은 학생이 \textbf{승리}합니다.
  \item 승리하면 베팅한 씨앗을 상금으로 받습니다.
  \item AI 대결도 가능합니다.
\end{enumerate}

\subsection{전적}
승리 수, 패배 수, 승률, 최고 연승이 기록됩니다.

\newpage
%% ═══════════════════════════════════════════
\section{테스트 창고}
%% ═══════════════════════════════════════════
\goalbox{시험을 응시하고, 성적을 분석하며, 오답을 복습합니다.}

\screenshot{08_tests.png}{0.85}

\subsection{시험 목록}
\begin{itemize}[nosep,leftmargin=*]
  \item 레벨별로 등록된 시험이 표시됩니다.
  \item 출제자별 필터: 국어농장 출제 / 소속 기관 출제.
  \item 각 시험의 정보: 출제자, 시행일, 응시 상태.
\end{itemize}

\subsection{OMR 답안 입력}
\begin{itemize}[nosep,leftmargin=*]
  \item 실제 시험지를 보면서 OMR처럼 답안을 입력합니다.
  \item 제출 후 \textbf{즉시 채점}되어 성적표가 표시됩니다.
\end{itemize}

\subsection{응시 이력}
\begin{itemize}[nosep,leftmargin=*]
  \item 응시한 모든 시험의 기록: 시험명, 점수, 정답률, 제출일.
  \item 통계: 응시 횟수, 평균 점수, 평균 정답률, 최고 점수.
  \item 점수 추이 차트로 실력 변화를 확인합니다.
\end{itemize}

\subsection{오답 노트}
특정 시험에서 틀린 문제만 모아 복습할 수 있습니다.

\newpage
%% ═══════════════════════════════════════════
\section{통합 성적표}
%% ═══════════════════════════════════════════
\goalbox{전 영역 학습 성과를 종합 분석합니다. 기간을 선택하여 특정 기간의 성과를 조회할 수 있습니다.}

\subsection{학습 캘린더}
선택한 기간의 학습 활동이 날짜별로 표시됩니다. 학습한 날은 색상으로 구분됩니다.

\screenshot{report_01_calendar.png}{0.9}

\newpage
\subsection{역량 레이더 차트}
10대 핵심 역량(어휘력, 독해력, 구조 독해력, 문장 독해력, 논리사고력, 문법능력, 어어 적용 능력, 배경지식, 문제 분석·전략 수립, 선택지 분석·전략 수립)이 레이더 차트로 시각화됩니다.

\screenshot{report_02_radar.png}{0.9}

\newpage
\subsection{역량별 상세 분석}
각 역량의 점수와 정답률이 막대 그래프와 테이블로 표시됩니다.

\screenshot{report_03_competency.png}{0.9}

\newpage
\subsection{영역별 분석}
국어 학습의 주요 영역별 성과가 정리됩니다.

\screenshot{report_04_area.png}{0.9}

\newpage
\subsection{활동별 상세}
테스트, 농장 모드, 일일 퀴즈/독해, 프로 모드 각각의 성과가 표시됩니다.

\screenshot{report_05_detail.png}{0.9}

\newpage
%% ═══════════════════════════════════════════
\section{랭킹}
%% ═══════════════════════════════════════════
\goalbox{다른 학생들과 학습 성과를 비교하며 학습 동기를 높입니다.}

\screenshot{04_ranking.png}{0.85}

\begin{itemize}[nosep,leftmargin=*]
  \item \textbf{시즌 랭킹}: 매월 갱신. 해당 시즌의 학습 점수 기준 순위.
  \item \textbf{누적 랭킹}: 전체 기간의 누적 점수 기준 순위.
  \item \textbf{필터}: 레벨별(같은 레벨끼리) 또는 전체 통합 비교.
  \item 순위, 이름, 점수가 테이블로 표시됩니다.
\end{itemize}

\newpage
%% ═══════════════════════════════════════════
\section{수확 장부와 씨앗 관리}
%% ═══════════════════════════════════════════

\subsection{수확 장부}
\goalbox{완료한 모든 학습 활동의 기록입니다.}

\screenshot{12_harvest.png}{0.85}

\begin{itemize}[nosep,leftmargin=*]
  \item 날짜, 학습명, 레벨, 영역, 수행시간, 진행률, 정답률, 획득 씨앗이 기록됩니다.
  \item 학습의 히스토리를 한눈에 확인할 수 있습니다.
\end{itemize}

\newpage
\subsection{씨앗 시스템}
\goalbox{국어농장의 보상 시스템입니다. 5종류의 씨앗을 모아 활용하세요.}

\screenshot{13_seeds.png}{0.85}

\begin{itemize}[nosep,leftmargin=*]
  \item \textbf{씨앗 종류}: 밀, 쌀, 옥수수, 포도, 사과 (5종).
  \item \textbf{획득 방법}: 일일 학습, 농장 학습, 대결 승리, 챕터 통과.
  \item \textbf{사용 방법}: 대결 베팅, 쇼핑몰 상품 구매.
  \item \textbf{교환}: 씨앗 종류 간 교환 가능.
  \item 씨앗 원장에서 모든 획득/사용 기록을 확인합니다.
\end{itemize}

\newpage
%% ═══════════════════════════════════════════
\section{커뮤니티}
%% ═══════════════════════════════════════════

\screenshot{09_community.png}{0.85}

5개 게시판이 있습니다:

\begin{tabularx}{\textwidth}{lXl}
\toprule
\textbf{게시판} & \textbf{설명} & \textbf{접근} \\
\midrule
학습 신청 게시판 & 학습 일정 신청 및 공유 & 구독 \\
커뮤니티 게시판 & 학습 이야기, 소식 공유 & 전체 \\
학습 질문 게시판 & 학습 관련 질문과 답변 & 구독 \\
학습 자료 게시판 & 관리자 제공 학습 자료 & 구독 \\
문의/상담 & 비공개 문의, 관리자 답변 & 전체 \\
\bottomrule
\end{tabularx}

\newpage
%% ═══════════════════════════════════════════
\section{쇼핑몰}
%% ═══════════════════════════════════════════

\screenshot{10_shop.png}{0.85}

교재와 교구를 구매할 수 있습니다.

\begin{itemize}[nosep,leftmargin=*]
  \item \textbf{교재}: 소쉬르 어휘 훈련 세트, 프레게 문법 집중 교재, 러셀 독해 전략 워크북, 비트겐슈타인 실전 모의고사.
  \item \textbf{교구}: 독해 전략 보드 게임, OMR 답안 입력 키트.
  \item 배송지를 등록하고 Toss Payments로 결제합니다.
\end{itemize}

\newpage
%% ═══════════════════════════════════════════
\section{과제 관리}
%% ═══════════════════════════════════════════
\goalbox{선생님이 배정한 과제를 확인하고 수행합니다.}

\begin{itemize}[nosep,leftmargin=*]
  \item \textbf{과제 유형}: 농장 학습 과제, 프로 학습 과제, 글쓰기 과제.
  \item \textbf{상태}: 진행중, 완료, 마감됨, 기한 초과.
  \item 과제를 클릭하면 해당 학습으로 바로 이동합니다.
  \item 학습을 완료하면 자동으로 과제 완료 처리됩니다.
  \item 학습 홈의 ``과제 바구니''에서 확인합니다.
\end{itemize}

\newpage
%% ═══════════════════════════════════════════
\section{구독 안내}
%% ═══════════════════════════════════════════

\screenshot{11_subscription.png}{0.7}

\subsection{무료로 이용 가능한 기능}
일일 퀴즈, 일일 독해, 대결, 역량 진단(최초 1회), 랭킹 확인, 커뮤니티 게시판(일부).

\subsection{구독 시 추가 이용 가능한 기능}
프로 모드, 농장별 모드(전체), 지식과 지혜(글쓰기), 테스트 창고, 학습 신청/질문/자료 게시판.

\vspace{0.5cm}
\begin{tabularx}{\textwidth}{Xlll}
\toprule
\textbf{기간} & \textbf{가격} & \textbf{할인} & \textbf{월 환산} \\
\midrule
1개월 & 65,000원 & -- & 65,000원 \\
3개월 & 175,500원 & 10\% & 58,500원 \\
6개월 & 312,000원 & 20\% & 52,000원 \\
12개월 & 546,000원 & 30\% & 45,500원 \\
\bottomrule
\end{tabularx}

\newpage
%% ═══════════════════════════════════════════
\section{내 정보}
%% ═══════════════════════════════════════════

\screenshot{14_profile.png}{0.7}

\begin{itemize}[nosep,leftmargin=*]
  \item \textbf{프로필 사진}: 업로드하여 변경할 수 있습니다.
  \item \textbf{이름}: 표시 이름을 수정합니다.
  \item \textbf{학년(레벨)}: 소쉬르1(초1) \textasciitilde{} 비트겐슈타인3(고3).
  \item \textbf{지역}: 17개 시도 중 선택.
  \item \textbf{학교}: 학교명을 입력합니다.
  \item \textbf{전화번호}: 학생 및 학부모 전화번호.
  \item \textbf{학습 시작일}: 달력 모드 또는 1일차 모드 선택.
  \item \textbf{비밀번호 변경}: 8자 이상의 새 비밀번호 설정.
\end{itemize}

\vfill
\begin{center}
\large\color{korfarm}\textbf{국어농장}\\[0.3cm]
\normalsize 스스로 완성하는 국어 근육\\
국어농장과 함께하세요.
\end{center}

\end{document}
